\section{Discussion}
\label{sec:discussion}

The development and empirical audit of GestureCAD provide several architectural, geometric, and human-factors insights into the integration of computer vision with geometric modeling.

\subsection{Architectural Decoupling and Web Performance}
A central design challenge in modern web-based graphics is managing the tension between declarative component state management (e.g., React, Vue) and high-frequency real-time rendering. Early iterations of web-based spatial interaction tools often routed tracking coordinates directly through component state variables. In our experiments, streaming 30--60 Hz landmark frames through React's Virtual DOM triggered continuous component re-renders, causing memory heap fragmentation and periodic garbage collection pauses ($\Delta t_{\text{gc}} > 25\text{ ms}$) that visibly disrupted drawing strokes.

GestureCAD's dual-path architecture resolves this tension:
\begin{enumerate}
    \item High-frequency spatial streams (cursor movements, live extrusion previews, plane manipulation) bypass the Virtual DOM entirely. Coordinates are maintained in mutable JavaScript references and written directly into Three.js \texttt{Float32Array} attribute buffers within the WebGL render loop (\texttt{useFrame}). This keeps cumulative software CAD execution under $0.2\text{ ms}$ per frame ($1.2\%$ of the $16.6\text{ ms}$ budget).
    \item Reactive state updates are reserved strictly for transactional boundaries (committing an extrusion, selecting a face, triggering Undo/Redo) managed via the Zustand store and Command Pattern.
\end{enumerate}
This architectural separation provides a practical blueprint for building high-performance spatial interaction systems within modern web frameworks.

\subsection{Topological Graph Validity vs. Spatial Geometric Embedding}
A key technical finding from our geometric audit is the vital distinction between \textit{topological graph invariance} and \textit{spatial geometric non-self-intersection}:
\begin{itemize}
    \item In Section~\ref{sec:geometry}, we showed that fan triangulation of planar caps yields a half-edge graph connectivity that strictly satisfies the Euler characteristic $\chi = V - E + F = 2$ and reports zero open boundary edges for any closed polygon.
    \item Furthermore, under Gauss's Divergence Theorem, overlapping exterior triangles contribute signed tetrahedral volumes of opposite sign that cancel algebraically, yielding exact agreement with analytical volume formulas even for non-convex profiles (as demonstrated in Benchmark B9).
    \item However, in physical $\mathbb{R}^3$ space, fan triangulation on concave profiles produces self-intersecting cap geometry. When exported to stereolithography (STL) files, self-intersecting facets violate the manifold assumptions of slicers and FEA meshers.
\end{itemize}
This insight underscores that topological validation ($\chi = 2$, $\partial \Omega = \emptyset$) is a necessary, but not sufficient, condition for geometric solid validity. Geometric modeling kernels must pair topological checks with explicit geometric non-self-intersection guarantees (e.g., via Ear-Clipping / Earcut triangulation).

\subsection{Practicality and Trade-Offs of Contactless CAD}
Contactless gesture CAD is not intended to replace high-precision desktop workstations for detailed engineering drafting. Rather, its practical utility emerges in specific contexts:
\begin{enumerate}
    \item \textbf{Sterile and Hazardous Environments}: In medical operating rooms (reviewing patient CT-scan geometry), pharmaceutical cleanrooms, or mechanical workshops with grease and debris, touching a mouse or keyboard contaminates the input peripheral. Contactless monocular webcam control provides an entirely hygienic interface.
    \item \textbf{Rapid Concept Ideation}: For industrial designers and architects, spatial hand gestures offer a fluid, natural medium for rough volumetric blocking and spatial proportion evaluation before committing to detailed drafting.
    \item \textbf{Ergonomic Diversity}: Sustained desktop mouse usage contributes to repetitive strain injuries (RSI) and carpal tunnel syndrome. GestureCAD provides an alternative physical modality for spatial manipulation.
\end{enumerate}
Conversely, the primary inherent trade-off of mid-air interaction is the \textbf{absence of haptic feedback}. When sketching on a physical tablet or using a mouse, physical surface friction stabilizes the human motor system. In mid-air, the lack of a physical backstop exacerbates physiological tremor. GestureCAD counterbalances this absence of tactile confirmation through software-assisted snapping, dead-zone gating, and clear visual state cues.

\subsection{Threats to Validity}
Several methodological factors must be considered when interpreting the reported results:
\begin{enumerate}
    \item \textbf{Synthetic vs. Real Noise Modeling}: While Benchmarks B3, B4, B10, and B11 evaluated comprehensive sweeps across 180 runs with multi-seed stochastic perturbations, the synthetic noise distributions (Gaussian perturbations and random-walk drift) may not fully capture the complex biomechanical dynamics of human motor tremor, fatigue-induced drift, or optical lighting flickers.
    \item \textbf{Hardware Dependency of Latency}: Software CAD latency (B1) was profiled on a high-performance host CPU (Intel Core i7-12700H). Performance may vary on lower-power client devices, mobile laptops, or under heavy browser background tab loads.
    \item \textbf{Unmeasured Motion-to-Photon Latency}: As discussed in Section~\ref{sec:experiments}, physical end-to-end latency remains unmeasured by high-speed optical hardware, representing an important gap between software execution time and perceptual responsiveness.
\end{enumerate}
